\scanpage{51}
Donc en admettant ce lemme, on trouve que $\forall$ syst. $v_1\ldots v_n$ peut se modifier galoisiennement de façon à donner lieu à un pt \deleted{générique} de $V(\widehat{\mathbf Z})$ \deleted{générique} tel que \uline{pour $\forall$ $\ell$}, $\forall\ell$ son image dans $V(\mathbf Z_\ell)$ soit générique rel. à $\mathbf Q$.

D'ailleurs, en sens inverse, si on suppose que $G$ satisfait à la condition d'approximation pour la place à l'infini,
\[
 G(\mathbf Q)\text{ dense dans }G(\widehat{\mathbf Q}),
\]
en ce cas \unsure{on voit que} $U.a$ contient un pt rationnel sur $\mathbf Q$. Donc on peut prouver que quitte à modifier les $v_i$ galoisiennement, on trouve des réseaux $M_{v_i}^{\mathbf Z}$ tous commensurables. [On fera cependant attention que le choix des $\sigma_i$ par lesquels on \inserted{doit} modifier $v_i$ \uline{dépend de $M$}; il est probable que pour $M$ variable, une place $v$ \unsure{est|pas} bien déterminée par la connaissance du foncteur fibre $\xi_v$].

\margin{Moyennant \inserted{\figure{\begin{tikzpicture}[baseline=(n.base)]\node[draw,ellipse,inner sep=2pt] (n) {le ``lemme'' et};\end{tikzpicture}}} la condition d'approximation sur $G$, la conjecture de Tate \unsure{renforcée} \uline{équivaut} à \deleted{\unsure{la}} l'énoncé d'indépendance des réseaux $M_{v_i}^{\mathbf Z}$ dans $M^{\mathbb A_f}$
\[
 (\mathbb A_f=\prod_{\ell\text{ fini}}\mathbf Z_\ell)
\]}
\mspage{50}
